\section{本轮结论与判定边界}
\label{sec:extension-summary}

\paragraph{衔接已有推导，而非重启旧问题。}
已有结果已经精确否定仅依赖 $\alpha$ 的一般候选界：一个二维、两更新实例满足 $\alpha=1$ 而 $\gamma=14/109<1/2$；任意固定正 $\alpha$ 的一般族可有 $\gamma\to0$。前轮已证明的全迹谱下界、共享模型精度价值上限和凸分配证书继续保留。本轮的工作是补上物理子类、近似余项、匹配系综与联合信息模型的范围，而不是把有限随机搜索改称证明。

\paragraph{物理子类的负结果已推进。}
新的三灯、两像素构造满足单位法向、Lambertian 着色、每灯对角场景 Jacobian、灯向切空间导数和正定标定先验；它在固定精化倍率下仍违反旧候选界，并有保持 $\alpha$ 固定的趋零边际比族。构造所用先验允许逐灯变化和通道耦合，因此并不自动解决共用同一对角先验等更窄子类。精确矩阵、物理参数与有理数核验在后文一并给出。

\paragraph{两项近似已有完整余项。}
对单位向量 $a$ 和指定的精确规范关系 $Aa=Bc$，本轮得到
\begin{equation}
 J_a(t)=\frac1{tq}+r+\delta_D+
            e^T(tI+S_\perp)^{-1}e,
 \qquad \delta_D\ge0.
\end{equation}
后两项分别量化空间信息不均匀和先验加权的规范残差。精确对齐只给出单侧近似，不保证误差小。固定二倍谱箱得到不以实测误差为输入的端点证书；两个端点共用谱质量，还可给价值差的带符号区间。direction-only 的大偏差由该区间识别，不再靠事后解释。完整证明保留英文数学正文，中文导读在本节及独立推导说明中给出。

\paragraph{匹配统计系综与历史幅值必须分开。}
同投影、同坐标不等于同统计系综。前轮修正后的 D 臂排序控制为 $R_A=0.55$，对象聚类区间为 $[-0.10,0.70]$；它替代旧的正向排序表述。经验幅值除以同坐标预测的中位数 $53.744$ 仍使用残差 bootstrap 基线，不能单独反驳匹配 GLS 定理。本轮 S1 先同时匹配噪声、估计器、基线、投影和模式坐标，再在 S2、S3 中逐步替换基线及有限扰动的前向模型。只有逐条件可比的比值才可作乘法归因，不能把中位数的乘积冒充总体恒等式。

\paragraph{谱界改进有明确而有限的紧度结论。}
在固定正定空间与未加权全迹下，留二候选的谱信息及完整集合条件数可加强原界。两维、两候选且固定基线条件数 $\chi$ 时，最坏下确界可精确刻画为 $1/\chi$，并给趋近参数族与收敛阶。该族允许更新尺度增长，因而不能把它解释成同时固定 $\kappa$、更新规模及条件数的联合最坏值。

\paragraph{排序桥接必须包含裕量。}
匹配 GLS 下，期望二次风险等于信息任务风险乘噪声尺度。协方差失配、偏差、线性化与采样误差给出可测扰动区间；只有风险差超过两侧误差区间之和的设计对，才得到排序保证。没有这一对间裕量，即使每项相对误差很小，排序也可反转。角度误差和尺度对齐误差不是无条件相同的二次任务，跨系统关系必须独立检验。

\paragraph{联合反照率与法向不破坏结构结论。}
联合模型在单位球切空间中是每像素三个自由度。Proper 标定先验可抬升尺度及更一般的分解规范；完全平坦先验下，完整非退化光度矩阵分解具有九维局部 $GL(3)$ 歧义，而不是只需删一个尺度。在固定商空间、固定线性化和固定任务下，$1-1/\kappa$ 上限与 Frank--Wolfe 证书仍成立；它们不是任意非线性重建的性能保证。
